% HD-S38. GPT-6 Astra Pro, /root/sixfolds; inherited effort not exposed.
% 21 November 2026. Actual normalized coefficient algebras; no global descent assumption.
\subsection{HD-S38: normalized coefficient algebras}
\label{sec:normalized-coefficient-algebras}

Keep the actual Gaussian source conventions of HD-S36 and HD-S37,
with a rank-nine period family of signature-\((3,3)\) sixfolds.
For coefficient ranks \(10\le m\le14\), every such source
whose actual group \(G\subset GL_m(K)\), \(K=\mathbb Q(i)\),
normalizes a proper nontrivial unital semisimple \(K\)-subalgebra
\[
           K\subsetneq\mathcal A\subsetneq\operatorname{End}_K(V)
                                                               \tag{S38.1}
\]
has ordinary Hermitian determinant \([-1]\) or \([-3]\).
The word \emph{normalizes} means the literal condition
\(\rho(g)\mathcal A\rho(g)^{-1}=\mathcal A\).
This includes all actual \(K\)-defined block systems and all
complex block systems defined over a finite Galois extension
whose projection algebra is Galois stable.
The latter qualification is not asserted for every complex
imprimitivity system.

The algebra in (S38.1) acts on the coefficient representation.
No action of that algebra on the primitive sixfold is asserted.
The actual abelian morphisms used in the proof arise only from
the induced-representation and quotient-curve constructions below.

\paragraph{1. Every proper \(K\)-block system in these ranks is covered.}
For an actual transitive \(K\)-block decomposition, choose one
block \(U\), its stabilizer \(H\), and
\(\sigma:H\to GL_K(U)\). The map
\[
  \bigoplus_{t\in G/H}U\longrightarrow V,
           \qquad (u_t)\longmapsto\sum_t\rho(t)u_t           \tag{S38.2}
\]
is the literal induced representation. Its complexification shows
that \(\sigma\) is absolutely irreducible: inducing a proper
complex invariant subspace of \(U\) would give a proper
\(G\)-invariant subspace of \(V\). The induced character
formula and the actual \(K\)-matrices give its exact character
field \(K\), as in (S36.2).
The blocks are orthogonal for the specified positive invariant
Hermitian coefficient form. Indeed start with a positive
block-diagonal \(K\)-Hermitian form and average it over \(G\).
The average remains block diagonal. Absolute irreducibility
makes this average and the specified invariant form positive
scalar multiples, so the latter has the same orthogonal blocks.

The entire quotient construction in paragraphs 2--3 of HD-S36
uses only (S38.2), not the origin of the blocks from reflections.
With \(H_0=\ker\sigma\), the actual curves are
\[
     Y=\widetilde C/H,\qquad C'=\widetilde C/H_0.
                                                               \tag{S38.3}
\]
The Shapiro function map, its inverse and the inversion of left
and right coset labels are exactly those now displayed in HD-S36.
For a vector \(v\in U(K)\), the normalized matrix-coefficient
idempotents satisfy the literal equality \(e_G=e_H\), because
the coefficient is zero off \(H\) and
\(m/|G|=\dim(U)/|H|\).
Thus pullback from \(JC'\) is an actual isogeny on the selected
primitive factor, with polarization multiplier \(|H_0|\) and
the same specified Gaussian corner. Its determinant multiplier
\(|H_0|^6\) is a norm. All these maps persist in the labelled
families after the same finite base change.

Here \(r=\dim_KU\le m/2\le7\). The rank-two bound is at
most seven; rank one is HD-S34. Rank three forces a sphere cover
with twelve reflection inertia, covered in HD-S31. For
\(4\le r\le7\), the parabolic formula
\(6=r(2g(Y)-2)+\sum c'_j\) and the marked-curve dimension
bound force \(g(Y)=0\). Rank nine then requires at least
twelve nontrivial branch values. Since
\(\sum c'_j=2r+6\le20<24\), the faithful block group has a
reflection. HD-S36 applies to it. This proves the asserted
determinant conclusion for every proper \(K\)-block system,
with the actual polarized maps retained.

\paragraph{2. Semisimple algebras reduce to blocks, tensors, or a field.}
If \(G\) contains a reflection, HD-S36 already gives the
conclusion, so suppose it is reflection-free. We may also suppose
that \(G\) is \(K\)-primitive, by paragraph 1.
The primitive central idempotents of \(\mathcal A\) give
\(K\)-subspaces of \(V\) permuted by \(G\). The action is
transitive by irreducibility. Consequently \(K\)-primitivity
makes the centre of \(\mathcal A\) a field \(E\).

First consider \(E=K\). Over \(\mathbb C\) the central
simple algebra becomes \(M_d(\mathbb C)\), acting on
\(V_{\mathbb C}\simeq\mathbb C^d\otimes\mathbb C^e\),
with \(de=m\). This identification is explicit: choose matrix
units \(a_{jk}\), a basis \(w_1,\ldots,w_e\) of
\(a_{11}V_{\mathbb C}\), and send
\(e_j\otimes f_s\) to \(a_{j1}w_s\).
Every normalizing \(g\) then acts as \(A_g\otimes B_g\).
Indeed its automorphism of the matrix-unit algebra is inner
(choose the image matrix-unit basis), and after multiplying by
the inverse implementer on the first factor it commutes with all
matrix units and acts only on the second factor.
Since \(\mathcal A\ne K\) and is proper in
\(\operatorname{End}_K(V)\), both \(d,e\) exceed one.
This is the actual fixed tensor-product decomposition excluded
in HD-S37, paragraph 2. Thus the remaining case has
\[
              K\subsetneq E\subset\operatorname{End}_K(V),
              \qquad [E:K]=a,\quad m=ar.                   \tag{S38.4}
\]
It suffices to use this centre; no splitting of the algebra over
\(E\) is needed.

\paragraph{3. The field-action quotient is forced to one endpoint.}
Conjugation gives an actual homomorphism
\[
               \gamma:G\longrightarrow\operatorname{Aut}_K(E).
                                                               \tag{S38.5}
\]
After complexification, the embeddings of \(E\) yield \(a\)
blocks of dimension \(r\). Their permutation by \(G\) is
transitive, since an orbit sum is invariant. An automorphism of a
field is determined by its composition with a fixed embedding;
therefore the automorphism action on embeddings is free. Its
transitivity implies that \(E/K\) is Galois and that (S38.5)
is surjective with regular image of order \(a\).

Let \(b_0\) be the number of inertia images nonidentity under
\(\gamma\). An image of order \(\ell\) has \(a/\ell\)
cycles on the embeddings. On a cycle, a fixed vector is determined
by one block vector and must satisfy the product holonomy.
Its fixed dimension is at most \(r\). Hence
\[
       c_j\ge m(1-1/\ell)\ge\lceil m/2\rceil
                   \qquad\text{when }\gamma(g_j)\ne1.       \tag{S38.6}
\]
If \(b_0\ge3\), then
\[
     \sum c_j\ge3\lceil m/2\rceil+2(12-3)>2m+6
                           \quad(10\le m\le14),             \tag{S38.7}
\]
a contradiction. Zero nonidentity images cannot generate the
nontrivial quotient, and one contradicts the product relation.
Thus exactly two images occur. They are inverse and generate
the full quotient, which is cyclic of order \(a\). Each has
codimension at least \(m-r\), so
\[
      2m+6\ge2(m-r)+2(b-2),\qquad12\le b\le r+5.           \tag{S38.8}
\]
Since \(a\ge2\) and \(m\le14\), this forces all equalities
\[
  m=14,\quad a=2,\quad r=7,\quad b=12,\quad b_0=2,
  \qquad(c_{\rm outside},c_{\rm inside})=(7,2).              \tag{S38.9}
\]

Write \(H=\ker\gamma\) and let \(t,u\) be the two outside
inertia elements. In the two complex blocks, either has form
\(\left(\begin{smallmatrix}0&A\\B&0\end{smallmatrix}\right)\).
Fixed dimension seven means \(AB=I\), hence \(t^2=u^2=1\).
The other ten inertia elements are \(E\)-linear and have
\(E\)-codimension one. Let \(N\) be their normal closure
in \(G\). It lies in \(H\). In \(G/N\) the product
relation identifies \(u=t^{-1}=t\), so \(|G/N|\le2\).
Since \(G/H\) has order two, \(N=H\).
Conjugation by an \(E/K\)-semilinear map preserves
\(E\)-rank; thus \(H\) is generated by actual
\(E\)-reflections on \(E^7\).

The representation on either complex block is irreducible.
A proper \(H\)-invariant subspace \(W\) in one block would
give the proper \(G\)-invariant subspace \(W\oplus tW\).
Its character cannot take all its values in \(K\). Otherwise
the two embedding representations of \(H\) are equivalent,
and \(V_{\mathbb C}\simeq U\otimes\mathbb C^2\) has
\(H\) acting on \(U\) alone. A fixed intertwiner implementing
the automorphism induced by \(t\) makes
\(t=A\otimes B\); since \(t^2=1\), rescale \(A,B\)
so that \(B^2=1\). An eigenline of \(B\) then yields an
actual proper \(G\)-invariant subspace, a contradiction.

\paragraph{4. The rank-seven reflection module gives actual rank-two blocks.}
Cohen's classification, Theorem (2.4), pp.386--387, real list
(4.4)(i), p.399, and Theorem (5.12), p.435 with table (4.16),
p.412, leaves in rank seven the rational-character real types
\(A_7,B_7,D_7,E_7\), and the monomial groups
\(G(d,e,7)\), \(e\mid d\).
The real types contradict the preceding character conclusion.
Use the actual group isomorphism to a standard monomial model
provided by the remaining classification case.

This model is defined over \(E\), as can be proved from literal
traces rather than assumed. Put \(\zeta=\exp(2\pi i/d)\).
The two elements
\[
 \operatorname{diag}(\zeta,\zeta^{-1},1,\ldots,1),\qquad
 \operatorname{diag}(\zeta,\zeta,\zeta^{-2},1,\ldots,1)
\]
belong to \(G(d,d,7)\subset H\). Their traces show that
\(s=\zeta+\zeta^{-1}\) and
\(v=2\zeta+\zeta^{-2}\) lie in \(E\). Since \(d>1\),
\(s\ne2\), and the actual identity
\[
                   \zeta=\frac{v-s^2+1}{2-s}                 \tag{S38.10}
\]
shows \(\zeta\in E\). With \([E:K]=2\) and
\(K=\mathbb Q(i)\), the cyclotomic degree bound gives
\(\operatorname{lcm}(4,d)\in\{4,8,12\}\).
The cases \(d=2,4\) again have character in \(K\), and
\(d=1\) is reducible in its seven-dimensional permutation
model. Thus
\[
                 d\in\{3,6,8,12\},\qquad E=K(\zeta).
                                                               \tag{S38.11}
\]
Both the given and standard modules have \(E\)-matrices.
Solving the finite intertwiner equations over \(E\), exactly
as in (G25.3), gives an invertible \(T\in GL_7(E)\).
This is a literal \(K\)-linear map on the fourteen-dimensional
space after restriction of scalars.

In the standard model write \(t=A\tau\), where
\(\tau\) is the nontrivial automorphism of \(E/K\) applied
entrywise. Explicitly the change by \(T\) replaces the old
semilinear coefficient matrix \(A_0\) by
\(A=T A_0\tau(T)^{-1}\).
The automorphism \(\tau\) sends \(\zeta\) to a coprime power and hence
preserves the actual set \(G(d,e,7)\). The matrix \(A\)
therefore normalizes that monomial reflection group linearly.
Its seven coordinate lines are preserved by every such normalizer.
The published uniqueness statement is Cohen (2.7), pp.389--390,
whose four exceptions all have rank at most four. The following
intrinsic description also specifies the lines being used.

When diagonal reflections have order greater than two, their root
lines are precisely the coordinate lines. When their order is
two, their conjugacy class has size seven; the phase-swap
reflection class has size \(21d\). Conjugating by diagonal
phases with compensation in a third coordinate makes that second
class transitive, so class size again distinguishes the axes.
In the remaining case \(e=d\), the root lines are
\(\mathbb C(e_j-\zeta^a e_k)\).
For \(d\ge4\), the graph joining lines with normalized
inner-product absolute value greater than \(1/2\) has one
connected component for every pair \(\{j,k\}\): roots on
different pairs have inner product zero or absolute value
\(1/2\), whereas consecutive phases on one pair join because
\(\cos(\pi/d)>1/2\).
For \(d=3\), the three root lines on one pair have identical
orthogonality neighborhoods, and different pairs have different
neighborhoods. Their common neighbors are exactly roots on
disjoint pairs; in rank seven this distinguishes the pairs.
The component spans, or these neighborhood-class spans, are the
coordinate pair-planes. Their nonzero pairwise intersections are
exactly the coordinate lines.

To apply this metric description to \(A\), observe that
\(A^*A\) commutes with the irreducible unitary group and is
therefore a positive scalar. Rescaling by its positive square
root makes \(A\) unitary without changing any line. Thus the
intrinsic description applies. Since \(\tau\) also fixes the
coordinate axes, both \(H\) and \(t\) permute the same seven
\(E\)-lines. These are seven actual \(K\)-subspaces of
dimension two, contradicting \(K\)-primitivity.

For a literal coordinate description, write
\(E=K(\sqrt q)\), \(q=2\) or \(3\). In the basis
\((1,\sqrt q)\), multiplication and field conjugation are
\[
 u+v\sqrt q\longmapsto
 \begin{pmatrix}u&qv\\v&u\end{pmatrix},\qquad
 \tau\longmapsto\begin{pmatrix}1&0\\0&-1\end{pmatrix}.
                                                               \tag{S38.12}
\]
Replacing each entry of \(T\) and \(A\) by the first block
gives the exact fourteen-dimensional \(K\)-maps. The seven
two-coordinate summands are the \(K\)-block system just
proved. Applying paragraph 1 would bound its period rank by
seven as well; no replacement of its Gaussian action is involved.

\paragraph{5. Exact residual scope.}
The centre-field case has been excluded, the centre-\(K\)
case is excluded by the actual tensor map, and product centres
give the already covered \(K\)-blocks. This proves the
statement for every actual algebra (S38.1).

For such a Galois-stable complex block system, write \(L/K\)
for a finite Galois field of definition and \(\mathcal D\)
for the \(L\)-span of its projection idempotents. Then
\(\mathcal D\) descends by coefficientwise fixed points to
the actual finite etale algebra
\(\mathcal A=\mathcal D^{\operatorname{Gal}(L/K)}\).
The descent map is \(\mathcal A\otimes_KL\to\mathcal D\),
\(a\otimes z\mapsto za\). An explicit inverse is supplied
by a basis \(b_j\) of \(L/K\) and its trace-dual basis
\(b_j^*\): for \(x\in\mathcal D\),
\[
 x=\sum_j b_j\operatorname{Tr}_{L/K}(b_j^*x),\qquad
 \operatorname{Tr}_{L/K}(b_j^*x)
   =\sum_{\gamma\in\operatorname{Gal}(L/K)}
                   \gamma(b_j^*)\gamma(x)\in\mathcal A.
\]
Conversely the embedding idempotents of such an algebra give the
system after complexification. These are the precise maps behind
the Galois-stable assertion at the beginning.

A remaining Gaussian source of determinant outside
\(\{[-1],[-3]\}\), coefficient rank ten through fourteen,
must therefore be complex imprimitive, \(K\)-primitive,
tensor indecomposable, and normalize no proper nontrivial
semisimple \(K\)-subalgebra. In particular none of its complex
block projection algebras can be Galois stable. No theorem
asserting the existence of a Galois-stable choice has been used.
The exact finite certificate checks the endpoint inequalities,
cyclotomic trace identity, field matrices and the root-line
reconstruction; the proof above supplies all group, coefficient,
curve and polarization maps.
